\setchoices{label=(\Alph*)}

\begin{question}[index=1]
设 $f(u)$ 可导，$z=xyf\!\left(\dfrac yx\right)$ 满足 $x\dfrac{\partial z}{\partial x}+y\dfrac{\partial z}{\partial y}=y^2(\ln y-\ln x)$，则 $f(u)$ 的极小值为
\begin{choices}[columns=2]
\item $\eu^{-1}$
\item $\dfrac12\eu^{-1}$
\item $-\dfrac12\eu^{-1}$
\item $\eu$
\end{choices}
\end{question}
\QuestionBreak

\begin{question}[index=2]
设反常积分 $\displaystyle I=\int_3^{+\infty}\frac{\dd x}{(x-1)(\ln x)^a}$（$a>0$）收敛，则 $a$ 的取值范围为
\begin{choices}[columns=4]
\item $(0,+\infty)$
\item $(\eu,+\infty)$
\item $(1,+\infty)$
\item $(1,\eu)$
\end{choices}
\end{question}
\QuestionBreak

\begin{question}[index=3]
设 $f(x)$ 满足 $f(0)=0$，则 $f'(0)$ 存在的充分必要条件是
\begin{choices}[columns=2]
\item $\displaystyle\lim_{x\to0}\frac{f(2x)-f(x)}x$ 存在
\item $\displaystyle\lim_{x\to0}\frac{f(x)-f(-x)}{2x}$ 存在
\item $\displaystyle\lim_{x\to0}\frac{f(\eu^x-1)}x$ 存在
\item $\displaystyle\lim_{x\to0}\frac{|f(x)|}x$ 存在
\end{choices}
\end{question}
\QuestionBreak

\begin{question}[index=4]
下列级数中条件收敛的个数为
\par\circlednumber{1} $\displaystyle\sum_{n=1}^\infty(-1)^n\left(1-\cos\frac1n+\frac1{\sqrt n}\right)$；
\par\circlednumber{2} $\displaystyle\sum_{n=2}^\infty(-1)^n\sin\frac1n$；
\par\circlednumber{3} $\displaystyle\sum_{n=1}^\infty\frac{(-1)^n}{n\sqrt[n]n}$.
\begin{choices}[columns=4]
\item $0$
\item $1$
\item $2$
\item $3$
\end{choices}
\end{question}
\QuestionBreak

\begin{question}[index=5]
设 $3$ 维列向量 $\vect\alpha_1,\vect\alpha_2,\vect\alpha_3$ 线性无关，$\vect A=(a_{ij})_{3\times3}$ 满足 $\vect A\vect\alpha_1=\vect\alpha_1$，$\vect A\vect\alpha_2=\vect\alpha_2-\vect\alpha_3$，$\vect A\vect\alpha_3=\vect\alpha_3-\vect\alpha_1$，$A_{ij}$ 是 $a_{ij}$ 的代数余子式，则 $A_{11}+A_{22}+A_{33}=$
\begin{choices}[columns=4]
\item $0$
\item $1$
\item $2$
\item $3$
\end{choices}
\end{question}
\QuestionBreak

\begin{question}[index=6]
\textcolor{red}{\bfseries 错题}\quad 设 $3$ 阶矩阵 $\vect A$ 满足 $\vect A^*=\vect A-2\vect E$，$|\vect A|=-1$，其中 $\vect A^*$ 为 $\vect A$ 的伴随矩阵，则下列命题中正确的个数为
\par\circlednumber{1} $\vect A-\vect E$ 相似于对角矩阵；
\par\circlednumber{2} $(\vect A-\vect E)^3=\vect O$；
\par\circlednumber{3} 方程组 $\vect A\vect X=\vect X$ 有两个线性无关的解向量；
\par\circlednumber{4} 方程组 $(\vect A-\vect E)^*\vect X=\vect0$ 有三个线性无关的解向量.
\begin{choices}[columns=4]
\item $1$
\item $2$
\item $3$
\item $4$
\end{choices}
\end{question}
\QuestionBreak

\begin{question}[index=7]
设 $\vect A,\vect B$ 均为 $n$ 阶矩阵，则 $\vect A$ 与 $\vect B$ 相似的充分条件是
\begin{choices}[columns=2]
\item $\vect A$ 与 $\vect B$ 有相同的特征值
\item $\vect A$ 与 $\vect B$ 有相同的特征向量
\item $\vect A^2$ 与 $\vect B^2$ 相似
\item $\vect A$ 与 $\vect B$ 均与同一对角矩阵相似
\end{choices}
\end{question}
\QuestionBreak

\begin{question}[index=8]
设二维随机变量 $(X,Y)$ 的概率密度为 $f(x,y)=\dfrac1{2\pi}\eu^{-\frac12(x^2+y^2)}$，$-\infty<x,y<+\infty$，则 $D(X^2+Y^2)=$
\begin{choices}[columns=4]
\item $1$
\item $4$
\item $2$
\item $\dfrac12$
\end{choices}
\end{question}
\QuestionBreak

\begin{question}[index=9]
设随机变量 $X_1,X_2,\ldots,X_n$（$n>1$）相互独立且服从同一分布，$D(X_i)=\sigma^2$（$i=1,2,\ldots,n$），$\sigma>0$，$\displaystyle\overline X=\frac1n\sum_{i=1}^nX_i$，$\displaystyle S^2=\frac1{n-1}\sum_{i=1}^n(X_i-\overline X)^2$，则下列选项正确的是
\begin{choices}[columns=2]
\item $S$ 与 $\overline X$ 相互独立
\item $S$ 是 $\sigma$ 的无偏估计量
\item $S$ 是 $\sigma$ 的最大似然估计量
\item $S$ 是 $\sigma$ 的一致估计量
\end{choices}
\end{question}
\QuestionBreak

\begin{question}[index=10]
设 $(X,Y)\sim N\!\left(1,0;9,16;-\dfrac12\right)$，$Z=\dfrac X3+\dfrac Y2$，则 $\rho_{XZ}=$
\begin{choices}[columns=4]
\item $0$
\item $1$
\item $2$
\item $3$
\end{choices}
\end{question}
